\section{Matchertext configuration variations}
\label{app:concrete:variations}

Even if adequately well-justified,
we cannot expect the standard matchertext configuration
of \cref{sec:design:concrete} to be a perfect or painless fit
for all situations in which string embedding is useful.
Including any matcher pair in $\Pi$ has the cost
of requiring those matchers to be escaped
in string literals and comments, for example,
as \cref{sec:embed} details.
There may be legitimate or even unavoidable reasons
to use other matchertext configurations in some cases,
keeping in mind that doing so reduces interlanguage embedding compatibility.

In general, deviations from the standard matchertext configuration
can be either \emph{tightening} (more restrictive),
\emph{loosening} (less restrictive), or a combination.

\subsection{Tightening variations}
\label{app:concrete:variations:tight}

A matchertext configuration $(\Sigma',\Pi')$
is a \emph{tightening} of
the standard matchertext configuration $(\Sigma,\Pi)$
if it only removes characters from the alphabet ($\Sigma' \subseteq \Sigma$)
and/or makes additional matcher pairs sensitive ($\Pi' \supseteq \Pi$).
A string in the tightened matchertext configuration
may be copied verbatim to an embedding context
expecting the standard matchertext configuration,
but not necessarily in the other direction.

As \cref{app:uri:liberal} details,
uniform resource identifiers (URIs)~\cite{rfc3986} traditionally allow
only graphical characters -- and no spaces or control codes for example --
in order to make them manually transcribable.
To serve this transcribability purpose,
significant spaces and control codes must not appear \emph{anywhere} in a URI,
even in an embedded matchertext substring.
Thus, URIs may represent a justifiable use-case
for an alternate matchertext configuration that removes
the non-graphical characters from the alphabet.
This would unfortunately mean that a string cannot, in general,
be copied from a standard matchertext language into a matchertext URI
without transformation (\ie escaping spaces and control codes).

\subsection{Loosening variations}

A matchertext configuration $(\Sigma',\Pi')$
is a \emph{loosening} of
the standard matchertext configuration $(\Sigma,\Pi)$
if it only adds characters to the alphabet ($\Sigma' \supseteq \Sigma$)
and/or removes sensitive matcher pairs ($\Pi' \subseteq \Pi$).
A string in the standard matchertext configuration
may be copied verbatim to an embedding context
supporting the loosened matchertext configuration,
but not necessarily in the other direction.

A loosened matchertext configuration might be justified,
for example, if it is deemed critical to embed strings in some language $L_e$
that frequently makes unmatched uses of some ASCII matchers,
and the pain of escaping or otherwise adapting that syntax is deemed too great.

Mathematical notation, for example,
sometimes uses
``mismatched'' parentheses and square brackets
to represent half-open/half-closed intervals.
That is,
$[0,1)$ typically means any real number $r$ greater than or equal to zero
but strictly less than one ($0 \le r < 1$).
A machine-readable language making frequent use of this mathematical notation
might be considered too painful to embed
in a standard matchertext configuration,
and therefore might ``demand'' a looser configuration
in which perhaps only the curly braces are sensitive as matcher pairs.

This example seems fairly hypothetical, however,
as extremely few machine-readable languages
appear to support this mathematical
half-open/half-closed interval notation anyway.
Languages that do support some form of half-open/half-closed syntax
often do so with other, more matchertext-friendly notation.
Swift~\cite{apple22swift}, for example,
supports \href{https://docs.swift.org/swift-book/LanguageGuide/BasicOperators.html#ID73}{\emph{half-open range} syntax} like \verb|1..<4|
for the sequence of integers starting from and including 1,
up to but not including 4.
This syntax is perfectly compatible with the standard matchertext configuration
because it uses the mathematical inequality operators,
rather than unmatched matchers,
to express the range's open upper endpoint.
